\begin{proposition}[5.7.5]
\label{IV.2.s5.7.5-fr}
Sean $X$ un preesquema localmente noetheriano y $\mathscr F$ un
$\mathcal O_X$-módulo coherente. Para que $\mathscr F$ verifique $(S_k)$
($k\geq1$), es necesario y suficiente que, para todo entero $r$ tal que
$0\leq r<k$, todo abierto $U$ de $X$ y toda sucesión
$(\mathscr F|U)$-regular $(f_i)_{1\leq i\leq r}$ de secciones de
$\mathcal O_X$ sobre $U$,
$(\mathscr F|U)/(\sum_{i=1}^r f_i(\mathscr F|U))$
carezca de ciclo primo asociado inmerso.
\end{proposition}

Demostremos primero la proposición para $k=1$; puede enunciarse también
diciendo que, \emph{para que $\mathscr F$ verifique $(S_1)$, es necesario y suficiente
que $\mathscr F$ carezca de ciclo primo asociado inmerso}. En efecto, decir que
$\mathscr F$ verifica $(S_1)$ significa que, en los puntos
$x\in\operatorname{Supp}(\mathscr F)$ tales que $\dim(\mathscr F_x)>0$, se
tiene $\operatorname{prof}(\mathscr F_x)\geq1$ (puesto que en los demás puntos
de $\operatorname{Supp}(\mathscr F)$ se tiene $\dim(\mathscr F_x)=0$, por
tanto $\operatorname{prof}(\mathscr F_x)=0$); pero decir que
$\operatorname{prof}(\mathscr F_x)\geq1$ significa que $\mathfrak m_x$
\emph{no está asociado a} $\mathscr F_x$ (\textbf 0, 16.4.6, (i)), o también que $x$
no está asociado a $\mathscr F$ (3.1.1); por otra parte, si $S$ es un
subpreesquema de $X$ cuyo espacio subyacente es
$\operatorname{Supp}(\mathscr F)$, se tiene (5.1.12.1)
$\dim(\mathscr F_x)=\dim(\mathcal O_{S,x})$; decir que
$\dim(\mathscr F_x)>0$ significa por tanto que $x$ no es un punto maximal de
$\operatorname{Supp}(\mathscr F)$ (5.1.2), de donde se sigue la conclusión.

En segundo lugar, mostremos que, para $k$ arbitrario $>1$, la condición del
enunciado es \emph{necesaria}. Es posible limitarse a considerar el caso en que $U=X$,
y nuestra afirmación quedará probada (por inducción sobre $k$ y en virtud de la
primera parte del razonamiento) si mostramos que, cuando $\mathscr F$ verifica
$(S_k)$ ($k>1$) y $f$ es una sección $\mathscr F$-\emph{regular} de
$\mathcal O_X$ sobre $X$, entonces $\mathscr F/f\mathscr F$ \emph{verifica}
$(S_{k-1})$ (es decir, verifica (5.7.2.1) donde $k$ se sustituye por $k-1$).
Ahora bien, para todo $x\in X$, $f_x$ es un elemento
$\mathscr F_x$-regular; si es invertible, se tiene
$\mathscr F_x/f_x\mathscr F_x=0$ y la conclusión es trivial. Si, por el
contrario, $f_x\in\mathfrak m_x$, se sabe que
$\operatorname{prof}(\mathscr F_x/f_x\mathscr F_x)
=\operatorname{prof}(\mathscr F_x)-1$ (\textbf 0, 16.4.6, (i)), y
$\dim(\mathscr F_x/f_x\mathscr F_x)=\dim(\mathscr F_x)-1$ (\textbf 0, 16.3.4), y
de ello resulta nuestra afirmación.

Probemos finalmente que, para $k>1$, la condición del enunciado es \emph{suficiente}.
Procederemos por inducción sobre $k$; sea $x$ un punto de $X$, y supongamos
primero que $\dim(\mathscr F_x)\geq k$. La hipótesis de inducción implica que
$\mathscr F$ verifica $(S_{k-1})$, por tanto
$\operatorname{prof}(\mathscr F_x)\geq k-1$; teniendo en cuenta (\textbf 0, 15.2.4),
existe, pues, una vecindad abierta $U$ de $x$ y una sucesión
$(\mathscr F|U)$-regular $(f_i)_{1\leq i\leq k-1}$ de secciones de
$\mathcal O_X$ sobre $U$, tales que $(f_i)_x\in\mathfrak m_x$ para
$1\leq i\leq k-1$. La hipótesis implica que
$\mathscr G=(\mathscr F|U)/(\sum_{i=1}^{k-1}
f_i(\mathscr F|U))$
carece de ciclo primo asociado inmerso; pero se tiene
$\dim(\mathscr G_x)=\dim(\mathscr F_x)-(k-1)\geq1$ (\textbf 0, 16.3.4), por tanto $x$
no está asociado a $\mathscr G$; se tiene, pues,
$\operatorname{prof}(\mathscr G_x)\geq1$ (\textbf 0, 16.4.6), y como
$\operatorname{prof}(\mathscr G_x)=\operatorname{prof}(\mathscr F_x)-(k-1)$
(\textbf 0, 16.4.6), se tiene $\operatorname{prof}(\mathscr F_x)\geq k$. Supongamos,
en segundo lugar, que $\dim(\mathscr F_x)=r<k$; como $\mathscr F$ verifica
$(S_{k-1})$, se tiene
$\operatorname{prof}(\mathscr F_x)\geq\inf(k-1,r)=r$, y esto termina la
demostración.

\begin{corollary}[5.7.6]
\label{IV.2.s5.7.6-fr}
Supongamos que $\mathscr F$ verifica $(S_k)$ ($k\geq1$); si
$(f_i)$ ($1\leq i\leq r$) es una sucesión $\mathscr F$-regular de secciones
de $\mathcal O_X$ sobre $X$ y si $r<k$,
$\mathscr F/(\sum_{i=1}^r f_i\mathscr F)$ verifica $(S_{k-r})$.
\end{corollary}

Esto resulta inmediatamente de (5.7.5).

\begin{corollary}[5.7.7]
\label{IV.2.s5.7.7-fr}
Para que $\mathscr F$ verifique $(S_2)$, es necesario y suficiente que
$\mathscr F$ carezca de ciclo primo asociado inmerso y que, para todo abierto
$U$ de $X$ y toda sección $(\mathscr F|U)$-regular $f$ de $\mathcal O_X$
sobre $U$, $(\mathscr F|U)/f(\mathscr F|U)$ carezca de ciclo primo asociado
inmerso.
\end{corollary}

\oldpage[IV-2]{106}
